\documentclass[12pt, letterpaper]{article}

\usepackage{../spreamble}

\title{CDA3101}
\author{Ian Zou}

\begin{document}

\maketitle

\part*{Von Neumann Architecture}

Instructions and data are stored in the same memory space. Single-memory, single-bus architecture.

\subsection*{Von Neumann Bottleneck}

Tasks are completed sequentially due to the single bus, and thus the speed of this single bus is a performance bottleneck.

\begin{itemize}
    \item \textbf{CU (Control Unit)}: Manages how the processor works by sending control signals, controlling input/output operations, and fetching instructions for execution.
    \item \textbf{ALU (Arithmetic Logic Unit)}: Performs arithmetic and logic calculations.
    \item \textbf{Memory Unit}: Consists of the main memory (RAM) and the registers. Single-memory system that holds both data and instructions.
    \item \textbf{Registers}: High-speed temporary storage on the CPU.
\end{itemize}

\part*{Measuring Performance}

\section*{Orders of Magnitude}
\begin{table}[h]
    \begin{center}
        \begin{tabular}{|p{2cm}|p{3cm}|}
            \hline
            \textbf{nano} & $ 10^{-9} $ \\
            \hline
            \textbf{micro} & $ 10^{-6} $ \\
            \hline
            \textbf{milli} & $ 10^{-3} $ \\
            \hline
            \textbf{kilo} & $ 10^3 $ \\
            \hline
            \textbf{mega} & $ 10^6 $ \\
            \hline
            \textbf{giga} & $ 10^9 $ \\
            \hline
            \textbf{tera} & $ 10^{12} $ \\
            \hline
        \end{tabular}
    \end{center}
\end{table}
\FloatBarrier

\section*{Units of Measure}
\begin{itemize}
    \item \textbf{Frequency}: \textit{Hertz}, or cycles per second. Used for \textbf{clock rate}.
    \item \textbf{Period}: Seconds per cycle. Used for \textbf{clock cycle}.
    \item \textbf{CPI}: Clocks per instruction. Highlights a CPU's \textbf{efficiency}.
\end{itemize}

Use \textit{dimensional analysis} to get to the desired unit.

\section*{Speedup}

$ \textrm{speedup} = \frac{t_{\textrm{old}}}{t_{\textrm{new}}} $

Where $ t_{\textrm{old}} $ and $ t_{\textrm{new}} $ are the times the program took before and after the upgrade, \textit{in seconds}.

\part*{discrete shit}

\section*{Two's Complement}

\textbf{Two's complement} is the standard for representing signed integers.

\begin{itemize}
\item \textbf{Most significant bit (MSB)} used as a sign bit.
\item For an $n$-bit word, the range of representable numbers is $ -2^{n-1} $ to $ 2^{n-1} - 1 $.
\end{itemize}

\subsection*{Two's Complement Representation}

\begin{enumerate}
    \item Take the binary of the absolute value of the integer.
    \item Flip the bits \textit{(one's complement)}.
    \item Add \texttt{1}.
\end{enumerate}

\part*{Assembly Structures}

\begin{itemize}
    \item Local variables stored on \textbf{stack} (either static or dynamic).
    \item Instructions and programs are stored in \textbf{text}.
\end{itemize}

\part*{Byte Order}

\begin{itemize}
    \item \textbf{Big-Endian (BE)}: Stores most significant byte first.
    \item \textbf{Little-Endian (LE)}: Stores most significant byte last.
\end{itemize}

\part*{RISC-V Instructions}

\begin{table}[h]
\begin{center}
    \begin{tabular}{|p{4cm}|p{6cm}|p{5cm}| }
        \hline
        \textbf{Mnemonic} & \textbf{Instruction} & \textbf{Description} \\
        \hline
        \texttt{add rd, rs1, rs2 } & Add & \texttt{rd <- rs1 + rs2}\\
        \hline
        \texttt{sub rd, rs1, rs2 } & Subtract & \texttt{rd <- rs1 - rs2}\\
        \hline
        \texttt{addi rd, rs1, imm12 } & Add (immediate r-value) & \texttt{rd <- rs1 + imm12}\\
        \hline
        \texttt{add rd, rs1, rs2 } & Add & \texttt{rd <- rs1 + rs2}\\
        \hline
    \end{tabular}
\end{center}
\end{table}

\FloatBarrier

\end{document}
